\documentclass[11pt]{article}
\usepackage{amsmath,amssymb,amsthm,geometry}
\geometry{margin=.9in}\setlength{\emergencystretch}{1.6em}
\newtheorem{theorem}{Theorem}\newcommand{\E}{\mathbb E}
\title{Quarter and half response from the initialized small hit}
\author{E5 continuation attempt}\date{Version 1}
\begin{document}\maketitle

\paragraph{Scope and notation.}
Fix the actual aligned balanced isotropic critical family
$S_0=cq^{1/\omega}$ with $c>0$, $\rho$ in the theory box,
$\lambda=\rho^2$, $\omega=(1-\lambda)/2$ and the WIT probe offset.
Time is normalized. Use physical Cartesian label vectors
$\Theta=(U,W)$, $dl=d\alpha/(2\pi)$, $|U|=|W|=\sqrt m$,
and population loss
$\mathcal L=\frac14\sum_{p=1}^2\E_{P_p}|r|^2$, $r(D)=D-f(D)$.
Let $Q=\lambda+q^2$ and $n=\mathsf m_2=\E_{P_2}[D_2f_2]/Q$.
The scalar $n$ below is an exact regression, not a radial logistic mass.
Use TBW's small-hit argument with its constants chosen smaller if
necessary; it gives a reached $\tau_0$ with $n(\tau_0)=\eta_0$
for a fixed sufficiently small $\eta_0>0$. The origin $t_h$
is the trained first right-half mass.

\section{The complementary residual and its signed exchange}
Set
\begin{equation}\label{an04rqtv1:defect}
 \mathcal R=4\mathcal L-Q(1-n)^2
  =\E_{P_1}|r|^2+
    \E_{P_2}\{r_1^2+[r_2-(1-n)D_2]^2\}\ge0.
\end{equation}
The equality is the exact Gaussian regression decomposition;
$\E_{P_2}D_2[r_2-(1-n)D_2]=0$.
Use the label Hilbert gradient and write
\[
 g=\nabla_\Theta n,\quad
 a=\tfrac Q2(1-n),\quad
 N=\dot\Theta-a g=-\tfrac14\nabla_\Theta\mathcal R,\quad
 \Gamma=\langle g,N\rangle .
\]
Indeed $\nabla\mathcal L=-\dot\Theta$, and differentiate
\eqref{an04rqtv1:defect}. In data coordinates
\[
 N=\tfrac12\E_{P_1}[J_D^T r]+
   \tfrac12\E_{P_2}[J_D^T(r-(1-n)D_2e_2)],
\]
where $J_D$ is the feature derivative with respect to both layers
at the actual population. Thus no source or residual component has
been discarded. The identities are, almost everywhere,
\begin{equation}\label{an04rqtv1:exchange}
 n'=a\|g\|^2+\Gamma,\qquad
 \mathcal R'=-4\|N\|^2-2Q(1-n)\Gamma .
\end{equation}
They use the actual finite-noise gradient flow and finite moments.

On the upper face $\mathcal R=r_*^2$ with $n<1/2$, the nonoutward
test is exactly
\begin{equation}\label{an04rqtv1:face}
 \Gamma\ge-\frac{2\|N\|^2}{Q(1-n)} .
\end{equation}
A strict version, or a locally Lipschitz scalar comparison that
supplies a uniform bound near the face, would exclude first exit.
No sign for this face is asserted for the critical trajectory.
The generic estimate $\Gamma\ge-\|g\|\|N\|$ does not give it:
apart from $N=0$, it suffices only where $\|N\|\ge a\|g\|$
(or additional direction information is known). Small complementary residual alone supplies
an upper norm bound, not this lower norm condition or signed angle.

\begin{theorem}[Initialized reached alternative]\label{an04rqtv1:alternative}
Fix $r_*=\lambda/800$. One can choose TBW's fixed $b,\eta_0$ so
that $\mathcal R(\tau_0)\le r_*^2/16$ for small $q$ and the same
compact Q2 arc reaches $\tau_0$ inside its strict vertical cone.
Let $\tau$ be the first time after $\tau_0$ with either $n=1/2$
or $\mathcal R=r_*^2$. For every fixed $k>8/\lambda$, small $q$ gives
\[
 \tau\le\tau_0+k\log(1/q).
\]
If the half hit occurs first, the actual trajectory has crossed
one-quarter response, and both response and cluster-loss criteria
are met at those crossing times for small $q$. If the residual
face occurs first, then
\begin{equation}\label{an04rqtv1:paidexit}
 -2Q\int_{\tau_0}^{\tau}(1-n)\Gamma\,dt
 =r_*^2-\mathcal R(\tau_0)+4\int_{\tau_0}^{\tau}\|N\|^2dt
 \ge15r_*^2/16 .
\end{equation}
Moreover $n(\tau)\ge\eta_0+[r_*^2-\mathcal R(\tau_0)]/(2Q)$.
This is a conclusion about the initialized path, with no
residual-cap hypothesis retained.
\end{theorem}
\begin{proof}
TBW's actual loss estimate before the small hit is
$\mathcal R\le(1-M(b))^2+\lambda(2\eta_0-\eta_0^2)+o_{q;b}(1)$.
Choose $b$ sufficiently large and then $\eta_0>0$ sufficiently
small that this is below $r_*^2/16$; use TBW's cone argument
with its tightened constants to reach this $\eta_0$.
At $t_h+b$ its compact Q2 arc has
$z=(U_2+W_2)/2\ge c_bq$, and
$|U_1|,|W_1|,|w|\le z/2$, $w=(U_2-W_2)/2$, with strict entry.
The cone is retained up to $\tau_0$ by the proved small-hit theorem.

Up to $\tau$, the full residual has bounded Gaussian $L^2$ norm
from loss decrease and \eqref{an04rqtv1:defect} bounds
$\|r\|_{P_1},\|r_1\|_{P_2}\le r_*$. Each of
$|R_{11}|,|R_{12}|,|R_{21}|$ is consequently at most $\lambda/100$
for small $q$, uniformly over receiver gates, by Cauchy--Schwarz
on the two clusters; the remaining $P_2$ contribution in $R_{21}$
costs $q\|r_2\|_{P_2}/2$. On an input cone with $U_2/|U|\ge1/\sqrt2$,
\[
 R_{22}\ge Q(1-n)/2-O(q)-O(e^{-d/q^2})\ge\lambda/5 .
\]
This subtracts the $P_2$ wrong gate by Cauchy--Schwarz and a
Gaussian tail; the $P_1$ term is at most $q r_*/2$ in absolute
value. It requires $n<1/2$, not monotonicity of $n$.

For the same compact arc the exact $z$ equation gives
$z'\ge\lambda z/8$ in the closed cone. At either face
$|U_1|=z/2$ or $|W_1|=z/2$, the upper speed is at most
$2(\lambda/100)z<z'/2$. At $|w|=z/2$,
$(|w|)'\le-\lambda z/10+\lambda z/200<z'/2$.
These strict signs include both signs of each first component,
both antisymmetric faces and every corner. The positive $z$ face
is protected by its growth inequality. The cone therefore persists
until the stopping event, even if the weak response temporarily
falls. Hence its amplitude is at least
$c_bq^2e^{\lambda(t-\tau_0)/4}$, while the full radial mass is
at most $C(1+t-t_h)$ from BRW. TBW bounds
$\tau_0-t_h=O(\log(1/q))$. At elapsed $k\log(1/q)$ with
$k>8/\lambda$, these estimates are inconsistent unless one of
the stopping events has already occurred.

At a first quarter or half hit while the cap holds,
$|\mathsf m_1-1|\le r_*/\sqrt{1+q^2}$ by regression on $P_1$,
$\mathcal L_1\le r_*^2/2$, and
$2\mathcal L_2\le Q(1-n)^2+r_*^2$.
The isotropic initial cluster losses tend to $1/2$ and
$\lambda/2$; these bounds satisfy the DRR thresholds
$1/2,1/4$ and loss ratio $169/196$ at both specified hits.
They do not align those times with a WIT learned window.
If instead the residual face is first, integration of
\eqref{an04rqtv1:exchange} proves \eqref{an04rqtv1:paidexit}.
Loss decrease gives
$r_*^2-\mathcal R(\tau_0)
 \le Q[(1-\eta_0)^2-(1-n(\tau))^2]$.
The positive difference implies $n(\tau)>\eta_0$ since
$n(\tau)<1/2$. Dividing by $2-\eta_0-n(\tau)\le2$ yields the
last bound. For simultaneous events the half-hit conclusions hold.
\end{proof}

\paragraph{Does restarting the loss budget pay the face?}
For any actual starting time $t_a$ with $n(t_a)=\eta_a$,
\begin{equation}\label{an04rqtv1:reset}
 \mathcal R(t)= \mathcal R(t_a)+
 Q[(1-\eta_a)^2-(1-n(t))^2]
 -4\int_{t_a}^t\|\dot\Theta\|^2\,du .
\end{equation}
This follows either from loss decrease as an equality or from
\eqref{an04rqtv1:exchange}. Restarting it at successive response
levels preserves all previous terms by telescoping.
To keep the same cap at a candidate quarter hit one needs the
actual dissipation to satisfy
\[
 4\int_{\tau_0}^{t_{1/4}}\|\dot\Theta\|^2dt
 \ge \mathcal R(\tau_0)+Q[(1-\eta_0)^2-9/16]-r_*^2;
\]
at half response replace $9/16$ by $1/4$.
These endpoint inequalities alone would not protect the preceding
cone; the intermediate version of \eqref{an04rqtv1:reset} is needed.
Cauchy--Schwarz using $n'=\langle g,\dot\Theta\rangle$ yields only
$\int\|\dot\Theta\|^2\ge(\eta-\eta_0)^2/\int\|g\|^2$.
It neither supplies the displayed lower dissipation nor signs
\eqref{an04rqtv1:face}. Thus subdividing into many small response
increments does not turn the original loss budget into a stronger
one. We cannot choose the half-hit alternative over the
residual-exit alternative on the actual critical trajectory.

\section{What is the unsigned transfer?}
The correlation at the face has a concrete full-kernel form.
For data vectors $D,D'$, write the exact two-layer tangent kernel
\[
 K_\Theta(D,D')=\int\left[
 (a^TD)_+(a^TD')_+ I_2+
 (D^TD'){\bf1}_{a^TD>0}{\bf1}_{a^TD'>0}bb^T\right]\,d\mu .
\]
Let $\tilde r_2(D')=r(D')-(1-n)D'_2 e_2$. Then
\begin{align}
 \Gamma={}&\frac1{2Q}\E_{D\sim P_2,D'\sim P_1}
       [D_2 e_2^TK_\Theta(D,D')r(D')]\nonumber\\
 &+\frac1{2Q}\E_{D,D'\sim P_2}
       [D_2 e_2^TK_\Theta(D,D')\tilde r_2(D')] .
 \label{an04rqtv1:coupling}
\end{align}
The Gaussian variables in each expectation are independent.
The weak residual has $\E_{P_2}D_2 e_2^T\tilde r_2=0$.
This removes a plain regression projection; the selected feature
kernel in \eqref{an04rqtv1:coupling} remains, so the identity
does not set the second term to zero. The $P_1$ term is also
present after the first concept meets a response threshold.
All expectations are on the actual initialized state. No sign
on \eqref{an04rqtv1:coupling} follows from positivity of the
data tangent kernel, which signs its quadratic forms with identical
data functions rather than this mixed correlation.

\section{Probe forcing if the response passage is reached}
Fix the deep COD witness $t_0=t_h+s_0$ before $\tau_0$ and retain
DRR's witness basis $u=r_{t_0}(\xi)/R_0$, $v=(-u_2,u_1)$ and
$R_0=|r_{t_0}(\xi)|$. In parameter space define
\[
 p_B=\nabla_\Theta[u^Tf(\xi)]/q^2,\qquad
 p_A=\nabla_\Theta[v^Tf(\xi)]/q .
\]
The current probe gates are part of these gradients. Their
integrals against $\dot\Theta$ give the DRR signed motions $B,A$.
For $g\ne0$ put
\[
 b_i=\langle p_i,g\rangle/\|g\|^2,\qquad h_i=p_i-b_i g
       \quad(i=B,A).
\]
On the zero set of $g$ put $b_i=0,h_i=p_i$. The identities below
are interpreted a.e.; no time derivative of these coefficients is
taken. All integrals can alternatively be stopped where
$\|g\|$ exceeds a cutoff and passed to the limit in their combined
exact expressions. In fact the terms are individually integrable
on finite time intervals: $|b_i n'|\le\|p_i\|\|\dot\Theta\|$
and $|\langle h_i,N\rangle|\le\|p_i\|\|\dot\Theta\|$
since $n'=\langle g,\dot\Theta\rangle$ and $h_i\perp g$.
Finite-moment feature bounds and dissipation pay these estimates.
The exact equations are
\begin{equation}\label{an04rqtv1:probe}
 B' = b_B n'+\langle h_B,N\rangle,\qquad
 A' = b_A n'+\langle h_A,N\rangle .
\end{equation}
They follow from $\dot\Theta=a g+N$ and
$n'=a\|g\|^2+\Gamma$, without any monotonicity of the response.
Thus the very residual update $N$ whose projection $\Gamma$
controls the response tube also has a probe-sensitive component
orthogonal to $g$. SPM's actual signed sources are included in
this update. Bounding $\Gamma$ alone does not bound the two
orthogonal correlations.

For any finite terminal time $T$, let
\[
 Z_i(T)=\int_{t_0}^T b_i n'\,dt+
                  \int_{t_0}^T\langle h_i,N\rangle\,dt .
\]
Thus $B=Z_B,A=Z_A$, with both populations evolved throughout.
For the actual learned candidate set $\mathcal T_q$ of DRR
(the half-response-centered fixed window with the response and
loss criteria), the remaining exact depth condition is
\begin{equation}\label{an04rqtv1:depth}
 \inf_{T\in\mathcal T_q,\ T\ge t_0}
 \left\{\frac{Z_A(T)^2}{2}+\frac{q^2Z_B(T)^2}{2}
                            -R_0Z_B(T)\right\}<e_*-e_q(t_0).
\end{equation}
An empty set gives no witness. Even proving the half-hit alternative
in Theorem~\ref{an04rqtv1:alternative} would only prove learning
at crossing times; aligning a suitable learned interval and
\eqref{an04rqtv1:depth} would still be required. In particular the
probe value at the moving small hit $\tau_0$ is not supplied by
COD's fixed-delay theorem.

The cap $\mathcal R\le r_*^2$ is an order-one data accuracy
estimate. It permits only
$\|N\|\le C\sqrt{S\mathcal R}$ by the feature derivative bound,
$S=\int m\,dl$. Used on the right sides of \eqref{an04rqtv1:probe},
it does not deliver the $q$ and $q^2$ signed precision. SPM proves,
on the actual pre-weak entry, large projected-gradient uncertainty
even after a response measurement is removed. No sign or scaled
joint bound for these orthogonal correlations is established here.
A hypothetical positive response slope is therefore insufficient
to bound the regenerated probe forcing.

\paragraph{Disposition.}
Two attempts were checked. Restarted loss control gives the exact
exchange identity but leaves \eqref{an04rqtv1:face} unsigned.
Projecting onto the actual weak response accounts for that
exchange but leaves the orthogonal probe forcing in
\eqref{an04rqtv1:probe} unpaid. The reached alternative is proved;
we have not proved quarter or half response for this initialization,
nor a deeper learned swing-by minimum. A violation of the proposed
residual tube is allowed by the analysis and would obstruct this
particular cone argument, rather than refuting the true learned
trajectory. No assertion of which alternative actually occurs is made.

\paragraph{Source ranges.}
TBW \texttt{an04tbwv1:seed, smallhit} supply this initialization's
compact Q2 entry and the actual fixed small response hit, with
constants chosen before small $q$. BRW's physical mass identity
is global after its reached half mass; its delay is used only
before that small hit. COD main and trained and SPM early are
fixed-resident-time results. P exactflow and learning are exact
finite-noise identities; its finite-moment Gaussian kernel and
the a.e.\ probe chain rule justify the gradient formulas here.
WIT and DRR specify the comparator and the measured target window,
not a critical-family arrival theorem. Both signs of all cone faces
and the signed residual-face test have been stated. All constants
may depend on fixed $c,\rho,\zeta$; no empirical or simulated
trajectory has been used.
\end{document}
